\clearpage
\VerdierPhysicalPageAnchor{0134}%
\begingroup
\small
\setlength{\emergencystretch}{1.5em}%
\setlength{\abovedisplayskip}{1pt plus 1pt minus 1pt}%
\setlength{\belowdisplayskip}{1pt plus 1pt minus 1pt}%

\VerdierSemanticAnchor{ch2-2-1-15}
\noindent\textbf{2.1.15.}\quad
\VerdierSemanticAnchor{ch2-2-1-15-associated-system}%
Soit $\mathcal D'$ une sous-catégorie triangulée pleine d'une catégorie
triangulée $\mathcal D$. Le système multiplicatif $S(\mathcal D')$ (2.1.7)
est appelé le \emph{système multiplicatif associé à la sous-catégorie
$\mathcal D'$}.

\VerdierSemanticAnchor{ch2-2-1-16}
\smallskip
\noindent\textbf{Proposition 2.1.16.}\quad
\begingroup\itshape
\VerdierSemanticAnchor{ch2-2-1-16-setup}%
Soit $F\colon\mathcal D_1\to\mathcal D_2$ un foncteur exact entre deux
catégories triangulées (1.1.3). Soit $S(F)$ l'ensemble des morphismes de
$\mathcal D_1$ qui sont transformés par $F$ en isomorphismes de
$\mathcal D_2$.
\VerdierSemanticAnchor{ch2-2-1-16-kernel-subcategory}%
Soit $\mathcal D'(F)$ la sous-catégorie pleine de $\mathcal D_1$ définie par
les objets de $\mathcal D_1$ qui sont transformés par $F$ en objets nuls de
$\mathcal D_2$.

\VerdierSemanticAnchor{ch2-2-1-16-a}
\par\noindent\hspace*{2em}a)\enspace La sous-catégorie $\mathcal D'(F)$ est
une sous-catégorie triangulée strictement pleine (tout objet de
$\mathcal D_1$ isomorphe à un objet de $\mathcal D'(F)$ est un objet de
$\mathcal D'(F)$) et saturée.

\VerdierSemanticAnchor{ch2-2-1-16-b}
\par\noindent\hspace*{2em}b)\enspace L'ensemble $S(F)$ est un système
multiplicatif saturé de $\mathcal D_1$.

\VerdierSemanticAnchor{ch2-2-1-16-c}
\par\noindent\hspace*{2em}c)\enspace L'ensemble $S(F)$ est le système
multiplicatif associé à la sous-catégorie $\mathcal D'(F)$ (2.1.15).
\par\endgroup

\VerdierSemanticAnchor{ch2-2-1-16-proof-omitted}
\noindent La preuve est facile et laissée au lecteur qui pourra démontrer de
même la proposition ci-après :

\VerdierSemanticAnchor{ch2-2-1-17}
\smallskip
\noindent\textbf{Proposition 2.1.17.}\quad
\begingroup\itshape
\VerdierSemanticAnchor{ch2-2-1-17-setup}%
Soit $H\colon\mathcal D\to\mathcal A$ un foncteur cohomologique d'une
catégorie triangulée dans une catégorie abélienne. Soit $S(H)$ l'ensemble des
morphismes de $\mathcal D$ qui sont transformés, ainsi que leurs translatés,
en des isomorphismes de $\mathcal A$.
\VerdierSemanticAnchor{ch2-2-1-17-kernel-subcategory}%
Soit $\mathcal D'(H)$ la sous-catégorie pleine de $\mathcal D$ définie par les
objets de $\mathcal D$ qui sont transformés, ainsi que leurs translatés, en
des objets nuls de $\mathcal A$.

\VerdierSemanticAnchor{ch2-2-1-17-a}
\par\noindent\hspace*{2em}a)\enspace La sous-catégorie $\mathcal D'(H)$ est
une sous-catégorie triangulée strictement pleine et saturée.

\VerdierSemanticAnchor{ch2-2-1-17-b}
\par\noindent\hspace*{2em}b)\enspace L'ensemble $S(H)$ est un système
multiplicatif saturé de $\mathcal D$.

\VerdierSemanticAnchor{ch2-2-1-17-c}
\par\noindent\hspace*{2em}c)\enspace L'ensemble $S(H)$ est le système
multiplicatif associé à la sous-catégorie $\mathcal D'(H)$ (2.1.15).
\par\endgroup

\VerdierSemanticAnchor{ch2-2-2}
\medskip
\noindent\textbf{2.2.\quad Construction de la catégorie localisée.}

\VerdierSemanticAnchor{ch2-2-2-1}
\smallskip
\noindent\textbf{2.2.1.}\quad
\VerdierSemanticAnchor{ch2-2-2-1-slice-categories}%
Soient $\mathcal C$ une catégorie (non nécessairement additive) et $S$ un
système multiplicatif de $\mathcal C$. Désignons, pour tout objet $X$ de
$\mathcal C$, par $S/X$ (resp. $X\backslash S$) la catégorie dont les objets
sont les morphismes de $S$ de but $X$ (resp. de source $X$)
\VerdierSemanticAnchor{ch2-2-2-1-morphism-diagrams}%
et dont les morphismes sont les diagrammes commutatifs :

\VerdierSemanticAnchor{ch2-2-2-1-s-over-x-diagram}
\begin{center}
  \begin{tikzpicture}[>=stealth,line width=.38pt,scale=.90,transform shape]
    \node (Y)  at (-1.55,1.05) {$Y$};
    \node (Yp) at ( 1.55,1.05) {$Y'$};
    \node (X)  at ( 0,-.95) {$X$};
    \draw[->] (Y) -- node[above=1pt] {$u$} (Yp);
    \draw[->] (Y) -- node[left=2pt] {$s$} (X);
    \draw[->] (Yp) -- node[right=2pt] {$t$} (X);
    \node at (3.15,.05) {$s,t\in S$};
  \end{tikzpicture}
  \quad ,
\end{center}

\VerdierSemanticAnchor{ch2-2-2-1-p135-boundary}
\endgroup
